%% heavyhex.tex — CppSim Heavy-Hex paper
%% Short paper: Tensor Network Simulation on Heavy-Hex Lattices using AMD GPUs
%%
%% Build: pdflatex heavyhex.tex

\documentclass[10pt,twocolumn]{article}
\usepackage[top=2cm,bottom=2cm,left=1.5cm,right=1.5cm,columnsep=0.5cm]{geometry}
\usepackage{graphicx}
\usepackage{amsmath,amssymb}
\usepackage{bm}
\usepackage{hyperref}
\usepackage{booktabs}
\usepackage{xcolor}
\usepackage{microtype}
\usepackage{abstract}
\renewcommand{\abstractnamefont}{\normalfont\bfseries}
\renewcommand{\abstracttextfont}{\normalfont\small}

%% TODO markers
\newcommand{\TODO}[1]{\textcolor{red}{\textbf{[TODO: #1]}}}
\newcommand{\CITE}[1]{\textcolor{blue}{[#1]}}
\newcommand{\chimax}{\chi_{\max}}

\begin{document}

\title{Tensor Network Simulation of the Heisenberg Model\\
       on Heavy-Hex Lattices using AMD GPUs}

\author{Paolo D. \\ \small AMD Research}

\date{\today}

% ─────────────────────────────────────────────────────────────────────────────
\maketitle

\begin{abstract}
We extend CppSim, a high-performance C++/HIP tensor network simulator,
to the heavy-hex lattice geometry used in IBM superconducting quantum processors.
The heavy-hex graph (degree-2 and degree-3 sites, 3-colorable gate schedule)
is implemented as a drop-in grid class with no changes to the gate application
or belief propagation kernels.
We simulate the isotropic Heisenberg model $H = J\sum_{\langle ij\rangle}
(\mathbf{X}_i\mathbf{X}_j + \mathbf{Y}_i\mathbf{Y}_j + \mathbf{Z}_i\mathbf{Z}_j)$
in the Heisenberg picture using Pauli transfer matrices,
and study the autocorrelation function $C(t) = \langle Z_c(t) Z_c(0) \rangle$
and operator lightcone on the \texttt{heavy\_hex\_3x3} graph (68 sites, 76 bonds).
We demonstrate chi convergence up to $\chi = $ \TODO{X} on an AMD MI100 GPU (32 GB),
and show a \TODO{X}$\times$ performance advantage of the C++/HIP implementation
over a Julia/AMDGPU.jl reference at saturated bond dimension.
\end{abstract}

% ─────────────────────────────────────────────────────────────────────────────
\section{Introduction}
\label{sec:intro}

Heavy-hex is the lattice topology used in IBM's superconducting quantum processors
\CITE{IBM heavy-hex}. Each qubit site has degree 2 or 3, and the graph admits
a natural 3-coloring of its bonds, enabling parallel gate application within
each color class.

Classical tensor network (TN) simulation of quantum circuits on heavy-hex graphs
serves two purposes: (i) validation of quantum computations for small system sizes
where exact simulation is feasible, and (ii) computation of noise-free reference
values for quantum error mitigation protocols such as the Multi-Product Formula
(MPF) \CITE{MPF ref}.

In this paper we show that CppSim \CITE{cppsim}, originally developed for
rectangular 2D Ising dynamics, extends naturally to heavy-hex topology with
minimal code changes. The hot-path kernels — gate application, belief propagation,
workspace management — are topology-agnostic and require no modification.
Only the grid construction and Hamiltonian gate matrix need to be defined
for the new geometry.

% ─────────────────────────────────────────────────────────────────────────────
\section{Heavy-Hex Lattice and Hamiltonian}
\label{sec:grid}

\subsection{Graph structure}

The \texttt{heavy\_hex\_$n_x \times n_y$} graph has:
\begin{itemize}
    \item $(n_x+1)(n_y+1)$ \emph{junction} sites (degree 3)
    \item remaining \emph{arm} sites (degree 2) connecting adjacent junctions
    \item $2n_x n_y + n_x + n_y$ bonds total
\end{itemize}

For $n_x = n_y = 3$: 16 junction sites, 52 arm sites, 76 bonds.
The graph is bipartite and 3-edge-colorable — bonds partition into 3 color
classes of sizes 25, 25, 26, corresponding to the natural gate layers
of the IBM hardware schedule.

Site tensors have shapes $[\chi^2, D]$ (arm) and $[\chi^3, D]$ (junction),
where $D=4$ is the Pauli basis dimension $\{I, X, Y, Z\}$.
The memory cost is dominated by junction sites:
$16 \times \chi^3 \times D \times 4$\,B (float32).

\subsection{Implementation}

The \texttt{FlatHeavyHexGrid} class inherits from \texttt{Grid} and populates
the same \texttt{Site}/\texttt{Bond}/\texttt{Gate} data structures used by the
rectangular grid. The color class assignment follows the hardware gate schedule
(JSON layers 0, 1, 2). No changes are required in \texttt{gate.hpp},
\texttt{measures.hpp}, or \texttt{workspace.hpp} beyond parametrizing the
workspace by \texttt{max\_degree=3} instead of 4.

\subsection{Hamiltonian and gate matrix}

We simulate the isotropic Heisenberg (XXX) model:
\begin{equation}
    H = J \sum_{\langle ij \rangle} \bigl(X_i X_j + Y_i Y_j + Z_i Z_j\bigr)
\end{equation}
with uniform coupling $J = 1$ and time step $\delta t = 0.1$.

In the Heisenberg picture with Pauli basis, the two-site gate
$U = e^{-i\theta(X\otimes X + Y\otimes Y + Z\otimes Z)}$
($\theta = J\delta t$) is represented by a $16\times 16$ real
Pauli transfer matrix (PTM) with block structure:
\begin{align}
    \{II, XX, YY, ZZ\} &\to \text{invariant (diagonal = 1)} \\
    \{IX, XI, YZ, ZY\} &\to \text{$4\times 4$ mixing block} \\
    \{IY, YI, XZ, ZX\} &\to \text{$4\times 4$ mixing block} \\
    \{IZ, ZI, XY, YX\} &\to \text{$4\times 4$ mixing block}
\end{align}
with elements $\alpha = \cos^2(2\theta)$, $\beta = \sin^2(2\theta)$,
$\gamma = \sin(2\theta)\cos(2\theta)$.
This was verified against a Julia reference implementation.

% ─────────────────────────────────────────────────────────────────────────────
\section{Autocorrelation and Chi Convergence}
\label{sec:ct}

\subsection{Autocorrelation function}

We initialize the Heisenberg-picture operator as $Z$ at a single center site
(site 13, a degree-3 junction at the graph midpoint) and $I$ elsewhere.
The autocorrelation $C(t) = \langle Z_c(t) Z_c(0) \rangle$ is computed
via the Bethe approximation after each Trotter step.

\subsection{Chi convergence}

\TODO{Insert ct\_convergence.png figure}

\TODO{Insert convergence table from plot\_ct\_convergence.py}

Figure~\ref{fig:ct} shows $C(t)$ for $\chi = 50, 100, 150, 200, 250, 300,
350, 400, 430$ over 30 Trotter steps. The result converges to
\TODO{X decimal places} at $\chi = $ \TODO{X}, indicating that
$\chi_{\max} = $ \TODO{X} is sufficient for faithful simulation on this graph.

The non-monotonic behaviour of $C(t)$ — a decrease followed by partial
recovery — is a genuine quantum recurrence of the XXX model on the degree-3
lattice, confirmed by its independence from $\chi$.

\subsection{Operator lightcone}

\TODO{Insert lightcone\_sites\_spatial\_site13.png figure (selected cycles)}

Figure~\ref{fig:lightcone} shows the non-identity Pauli weight $n(v,t)$
at each site for selected cycles. The operator spreads radially outward
from the center junction, following the 3-fold symmetry of the hexagonal
lattice. The wavefront reaches the boundary sites by cycle \TODO{X},
defining the computational window within which the simulation is exact
regardless of $\chi$ truncation at the boundary.

% ─────────────────────────────────────────────────────────────────────────────
\section{Performance}
\label{sec:perf}

\subsection{Chi scaling}

\TODO{Insert timing\_vs\_chi.png and chi\_growth.png}

Table~\ref{tab:perf} shows the wall-clock time per Trotter cycle at
saturated bond dimension for each $\chi$ value.
The dominant cost is the Householder QR decomposition
($\mathcal{O}(m n^2)$ where $m = \chi^2 D$ for degree-3 sites),
which exceeds the SVD cost above $\chi \approx $ \TODO{X} cycles.

\TODO{Insert chi scaling table}

\subsection{CppSim vs Julia/AMDGPU.jl}

\TODO{Insert comparison table (10x speedup result)}

At saturated bond dimension ($\chi = 430$, cycles 8--10), CppSim achieves
a $\sim$10$\times$ speedup over the Julia/AMDGPU.jl reference implementation
on the same AMD MI100 GPU (32\,GB HBM2). Both use Householder QR and
float32 arithmetic. The speedup arises from:
(i) zero GPU memory allocation during simulation (pre-allocated workspace),
(ii) direct HIP/rocBLAS calls without framework dispatch overhead, and
(iii) on-device \texttt{pseudo\_sqrt} kernel eliminating CPU--GPU transfers.

% ─────────────────────────────────────────────────────────────────────────────
\section{Conclusion}
\label{sec:conclusion}

We have demonstrated that CppSim, a C++/HIP tensor network simulator,
extends to the heavy-hex lattice with minimal code changes.
The topology-agnostic gate and belief propagation kernels accept any
\texttt{Grid}-derived class, and the heavy-hex grid requires only a new
site/bond/color construction.

On an AMD MI100 (GPU-B, 32\,GB), we achieve \TODO{exact} simulation up to
$\chi = $ \TODO{X} on the $3\times3$ heavy-hex graph (68 sites),
and demonstrate $\sim$10$\times$ speedup over a Julia reference implementation
at saturated bond dimension.

The operator lightcone visualization — computed as a byproduct of the
simulation with no additional cost — provides an intuitive picture of
operator spreading on the hexagonal lattice and defines the computational
window for quantum error mitigation protocols.

Future work will extend to larger graphs ($2\times2$, $4\times4$) on
higher-memory hardware (MI355X, 309\,GB), and implement interleaved
belief propagation to match Julia's online BP approach and eliminate
BP trap effects.

% ─────────────────────────────────────────────────────────────────────────────
\section*{Acknowledgments}
\TODO{Acknowledgments}

\bibliography{heavyhex}

\end{document}
